\section{Overview}

This section defines the scope of protection for the RCO protocol. The claims are structured to capture the fundamental principles of the invention while also covering specific implementations and variations. The objective is to establish broad protection over the core concept of feedback-coupled telemetry integrity, while ensuring that the claims remain precise and defensible.

\section{Independent Claims}

\subsection{Claim 1: Core Method}

A computer-implemented method for enforcing integrity of telemetry in a feedback-driven system, comprising:

receiving telemetry data from one or more sources;

transforming the telemetry data into a deterministic representation using a canonicalization process;

computing a cryptographic hash of the deterministic representation;

generating a lineage value by combining the cryptographic hash with a prior lineage state;

binding the lineage value to a system state representation to form a verification message;

distributing the verification message to a plurality of independent validator nodes;

obtaining, from at least a threshold number of the validator nodes, cryptographic signatures over the verification message;

aggregating the cryptographic signatures into a composite proof;

verifying the composite proof against a public verification key;

and accepting the telemetry data into the system only if the composite proof is valid,

wherein the method ensures that only telemetry data satisfying deterministic representation, lineage consistency, and state binding constraints is accepted.

\subsection{Claim 2: System Claim}

A distributed system for enforcing telemetry integrity, comprising:

an ingestion component configured to receive telemetry data and perform deterministic canonicalization;

a hashing component configured to compute cryptographic hashes of canonicalized data;

a lineage engine configured to generate chained lineage values based on prior lineage state;

a state binding component configured to combine lineage values with system state representations;

a plurality of validator nodes configured to independently verify the verification message and generate cryptographic signatures;

an aggregation component configured to combine signatures from a threshold number of validator nodes;

and a verification component configured to validate the aggregated signature and determine acceptance of telemetry data,

wherein the system enforces consistency across distributed nodes and prevents acceptance of invalid or inconsistent telemetry.

\subsection{Claim 3: Non-Transitory Medium}

A non-transitory computer-readable medium storing instructions that, when executed by one or more processors, cause the processors to perform the method of Claim 1.

\section{Independent Claim: Conceptual Abstraction}

\subsection{Claim 4: Feedback-Coupled Integrity Framework}

A method of maintaining correctness in a feedback-driven computational system, comprising:

treating telemetry as a causative input to system state evolution;

enforcing deterministic representation of telemetry;

establishing temporal lineage linking telemetry events;

binding telemetry to system state to preserve semantic consistency;

and requiring distributed verification of telemetry prior to state transition,

wherein the system prevents propagation of inconsistent or adversarial inputs through enforced verification at ingestion.

\section{Dependent Claims}

\subsection{Claim 5}

The method of Claim 1, wherein the deterministic representation comprises ordering telemetry fields according to a predefined canonical schema and encoding values into a normalized format to eliminate representational ambiguity.

\subsection{Claim 6}

The method of Claim 1, wherein the cryptographic hash function is selected from a class of collision-resistant hash functions, including but not limited to SHA-2, SHA-3, or equivalent constructions.

\subsection{Claim 7}

The method of Claim 1, wherein the lineage value is computed as a chained hash incorporating both the current telemetry hash and a prior lineage state to enforce temporal dependency.

\subsection{Claim 8}

The method of Claim 1, wherein the state binding includes a cryptographic representation of system state such that telemetry is verifiably associated with the state from which it was generated.

\subsection{Claim 9}

The method of Claim 1, wherein validator nodes operate independently and perform verification without reliance on a centralized authority.

\subsection{Claim 10}

The method of Claim 1, wherein the cryptographic signatures are generated using a threshold signature scheme.

\subsection{Claim 11}

The method of Claim 10, wherein the threshold signature scheme supports aggregation of partial signatures into a single compact proof.

\subsection{Claim 12}

The method of Claim 1, wherein telemetry data is rejected if the number of valid signatures does not meet a predefined threshold condition.

\subsection{Claim 13}

The method of Claim 1, wherein canonicalization, hashing, and lineage computation are performed locally at the telemetry source prior to distribution.

\subsection{Claim 14}

The method of Claim 1, wherein aggregation of signatures is performed by a dedicated aggregation service or distributed among validator nodes.

\subsection{Claim 15}

The method of Claim 1, wherein verification of the aggregated signature is performed by one or more independent verifier components.

\subsection{Claim 16}

The method of Claim 1, wherein the system supports batch processing of multiple telemetry events to improve throughput.

\subsection{Claim 17}

The method of Claim 1, wherein telemetry data is compressed or reduced to its hash representation prior to transmission to minimize network overhead.

\subsection{Claim 18}

The method of Claim 1, wherein validator nodes are dynamically added or removed without compromising the integrity guarantees of the system.

\subsection{Claim 19}

The method of Claim 1, wherein validator keys are periodically rotated to reduce long-term security risks.

\subsection{Claim 20}

The method of Claim 1, wherein the system operates across heterogeneous infrastructure environments including cloud, edge, and hybrid deployments.

\subsection{Claim 21}

The method of Claim 1, wherein the system enforces rejection of telemetry data under uncertain conditions to preserve integrity.

\subsection{Claim 22}

The method of Claim 1, wherein lineage states are periodically checkpointed to optimize storage and retrieval.

\subsection{Claim 23}

The method of Claim 1, wherein telemetry processing is parallelized across multiple validator nodes to improve scalability.

\subsection{Claim 24}

The system of Claim 2, wherein validator nodes are geographically distributed to improve fault tolerance and resilience.

\subsection{Claim 25}

The system of Claim 2, wherein the aggregation component supports asynchronous collection of signatures.

\subsection{Claim 26}

The system of Claim 2, wherein the verification component rejects telemetry data if any inconsistency is detected in lineage or state binding.

\subsection{Claim 27}

The system of Claim 2, wherein the storage component maintains an append-only log of accepted telemetry and associated proofs.

\subsection{Claim 28}

The method of Claim 1, wherein the system detects and rejects replayed telemetry by verifying lineage continuity.

\subsection{Claim 29}

The method of Claim 1, wherein adversarial inputs are detected through inconsistencies in deterministic representation or cryptographic verification.

\subsection{Claim 30}

The method of Claim 1, wherein the system maintains consistency across distributed nodes by ensuring identical lineage computation.

\subsection{Claim 31}

The method of Claim 1, wherein failure of individual validator nodes does not compromise system correctness as long as the threshold condition is satisfied.

\subsection{Claim 32}

The method of Claim 1, wherein the system degrades gracefully by rejecting telemetry under insufficient validation conditions while preserving prior state integrity.

\subsection{Claim 33}

The method of Claim 1, wherein telemetry is treated as a causative input influencing system state transitions.

\subsection{Claim 34}

The method of Claim 1, wherein deterministic processing reduces ambiguity and improves verifiability of telemetry data.

\subsection{Claim 35}

The method of Claim 1, wherein the system supports integration as a middleware layer between telemetry producers and consumers.

\subsection{Claim 36}

The method of Claim 1, wherein verification is performed in real time or near real time.

\subsection{Claim 37}

The method of Claim 1, wherein the system supports replay and reconstruction of system state using lineage data.

\subsection{Claim 38}

The method of Claim 1, wherein telemetry integrity is enforced prior to any downstream processing.

\subsection{Claim 39}

The method of Claim 1, wherein the protocol is implemented using standard cryptographic primitives without requiring specialized hardware.

\subsection{Claim 40}

The method of Claim 1, wherein the system is configured to operate under adversarial conditions while maintaining correctness guarantees.

\section{Novelty and Non-Obviousness Justification}

\subsection{Overview}

This section establishes the novelty and non-obviousness of the RCO protocol by systematically distinguishing it from existing systems and demonstrating that the invention represents a fundamental departure from prior approaches. The analysis addresses both structural and conceptual differences, and explains why the combination of features introduced by the protocol is not an obvious extension of existing techniques.

\subsection{Limitations of Existing Systems}

Existing systems for handling telemetry and data integrity fall into several broad categories, including logging systems, monitoring frameworks, cryptographic signing mechanisms, and distributed consensus systems.

Logging systems record telemetry for later analysis but do not enforce correctness at the point of ingestion. Monitoring systems detect anomalies but operate reactively, after data has already been accepted into the system. Cryptographic signing mechanisms provide authenticity but do not enforce deterministic representation or contextual binding. Distributed consensus systems ensure agreement among nodes but do not guarantee that the agreed-upon data is intrinsically valid or contextually consistent.

In all of these systems, telemetry is treated as an input that may be validated, analyzed, or corrected after it has been incorporated into system workflows. None of these approaches enforce a unified framework in which telemetry must satisfy multiple layers of constraints prior to acceptance.

\subsection{Absence of Feedback-Coupled Integrity}

A critical gap in prior art is the absence of feedback-coupled integrity enforcement. In existing systems, telemetry is not treated as a causative variable that directly influences system state in a formally constrained manner.

While feedback loops are common in control systems and distributed applications, there is no mechanism in prior art that ensures that telemetry entering such loops is simultaneously deterministic, cryptographically verifiable, temporally consistent, and contextually bound.

The RCO protocol introduces this coupling explicitly. Telemetry is not merely validated; it is transformed into a structured entity whose acceptance is contingent on satisfying a set of interdependent constraints. This creates a fundamentally different system behavior in which invalid telemetry cannot propagate into system state.

\subsection{Integration of Independent Mechanisms}

Individually, the components of the RCO protocol—canonicalization, hashing, lineage construction, state binding, and distributed signature verification—are known techniques. However, prior art does not disclose or suggest their integration into a single, unified pipeline that operates at ingestion time.

In existing systems, these mechanisms are used in isolation or in loosely coupled configurations. For example, hashing may be used for integrity checks, and signatures may be used for authentication, but there is no requirement that these mechanisms be applied sequentially and interdependently to enforce a global invariant.

The RCO protocol requires that all of these mechanisms operate together, with each stage reinforcing the constraints of the others. This interdependence is essential to the protocol’s ability to enforce deterministic, verifiable, and context-aware telemetry integrity.

\subsection{Non-Obviousness of the Combination}

The combination of features introduced by the RCO protocol is not an obvious extension of prior art for several reasons.

First, existing systems prioritize performance and flexibility over strict correctness at ingestion. Introducing deterministic canonicalization, lineage enforcement, and distributed verification at this stage would be considered unnecessarily restrictive in traditional designs.

Second, the idea of binding telemetry to system state is not suggested by existing approaches. Most systems treat telemetry and state as separate concerns, and there is no motivation in prior art to enforce a cryptographic relationship between them.

Third, the use of threshold-based distributed verification for telemetry ingestion is a departure from its typical use in consensus and blockchain systems. In those systems, signatures are used to validate transactions or blocks, not to enforce deterministic and context-aware telemetry integrity.

Finally, the protocol’s emphasis on rejecting uncertain data rather than attempting recovery or correction represents a shift in design philosophy that is not reflected in prior systems.

\subsection{Synergistic Effect}

The combination of the protocol’s components produces a synergistic effect that is greater than the sum of its parts.

Deterministic representation ensures consistent encoding. Cryptographic hashing provides integrity guarantees. Lineage enforces temporal consistency. State binding preserves semantic context. Distributed verification ensures trust without central authority.

Individually, these mechanisms address specific aspects of data integrity. Together, they create a system in which telemetry must satisfy all constraints simultaneously, resulting in a level of assurance that cannot be achieved by any single mechanism alone.

This synergy is not suggested by prior art and represents a key aspect of the invention’s novelty.

\subsection{Distinction from Blockchain Systems}

Blockchain systems are often cited as prior art for distributed verification and chained data structures. However, the RCO protocol differs fundamentally in purpose, structure, and operation.

Blockchain systems focus on achieving consensus over a sequence of transactions and ensuring immutability after inclusion. They do not enforce deterministic representation of input data, nor do they bind data to system state in a feedback-driven context.

Furthermore, blockchain validation occurs at the level of transactions and blocks, not at the level of telemetry ingestion within a feedback loop. The RCO protocol operates at a different layer, enforcing correctness before data enters the system rather than after it has been recorded.

These differences establish that the RCO protocol is not a variation of blockchain technology, but a distinct approach to data integrity.

\subsection{Distinction from Traditional Validation Systems}

Traditional validation systems perform checks on input data, but these checks are typically limited to schema validation, range checks, or rule-based filtering.

The RCO protocol goes beyond these approaches by introducing cryptographic and structural constraints that cannot be bypassed or overridden. Validation is not a separate step but an integral part of the data ingestion process.

This integration ensures that invalid data cannot enter the system, rather than relying on downstream mechanisms to detect and correct errors.

\subsection{Unexpected Advantages}

The RCO protocol provides several advantages that would not be expected from a straightforward combination of known techniques.

These include:

\begin{itemize}
\item elimination of downstream inconsistencies,
\item improved system reliability under adversarial conditions,
\item simplified debugging and auditing through verifiable lineage,
\item enhanced scalability due to deterministic processing.
\end{itemize}

These advantages arise from the interaction of the protocol’s components and are not achievable through conventional approaches.

\subsection{Generalization and Future Applicability}

The principles underlying the RCO protocol are applicable to a wide range of systems beyond the specific implementations described in this document.

This generality further supports the argument that the invention represents a new class of systems rather than a specific configuration of existing technologies.

\subsection{Conclusion}

The RCO protocol introduces a novel and non-obvious approach to enforcing telemetry integrity in feedback-driven systems. By integrating deterministic representation, cryptographic lineage, state binding, and distributed verification into a unified pipeline, the protocol addresses a problem that is not solved by existing systems.

The combination of these features, and the resulting system behavior, are not suggested by prior art and would not be obvious to a person skilled in the field. This establishes the novelty and non-obviousness of the invention and supports the claims presented in this document.